% !TEX program = tectonic
% =============================================================================
%  AULA 7 — Redes sem Fio: 5G e Redes Móveis de Nova Geração
%  Curso: Redes sem Fio  |  Profa. Anelise Munaretto
%  Conteúdo atualizado 2026: 5G SA/NSA, network slicing, MEC, mMIMO, mmWave.
%  Compilar: tectonic aula7.tex
% =============================================================================
\documentclass[aspectratio=169,11pt]{beamer}
% =============================================================================
%  PREÁMBULO COMÚN para as aulas de Redes sem Fio (Beamer + TikZ)
%  Incluido com % =============================================================================
%  PREÁMBULO COMÚN para as aulas de Redes sem Fio (Beamer + TikZ)
%  Incluido com % =============================================================================
%  PREÁMBULO COMÚN para as aulas de Redes sem Fio (Beamer + TikZ)
%  Incluido com \input{preamble.tex} en cada aula.
%  Paleta de cores e estilos são definidos aqui uma única vez.
% =============================================================================

\usepackage[T1]{fontenc}
\usepackage{amsmath,amssymb,mathtools}
\usepackage{graphicx}
\usepackage{tikz}
\usetikzlibrary{positioning,arrows.meta,calc,backgrounds,decorations.pathmorphing,fit,shadows,shapes.symbols,shapes.geometric}
\usepackage{pgfplots}
\pgfplotsset{compat=1.16}
\usepackage{tabularx,booktabs,multirow,array}
\usepackage[normalem]{ulem}
\usepackage{hyperref}

% ---------- Paleta de cores própria ----------
\definecolor{aneliseAzul}{RGB}{20,60,110}
\definecolor{aneliseVerde}{RGB}{30,140,70}
\definecolor{aneliseLaranja}{RGB}{215,120,20}
\definecolor{aneliseGris}{RGB}{90,90,90}

% ---------- Tema ----------
\usetheme{Madrid}
\usecolortheme{whale}
\setbeamercolor{structure}{fg=aneliseAzul}
\setbeamercolor{title}{fg=aneliseAzul}
\setbeamercolor{frametitle}{fg=aneliseAzul}
\setbeamercolor{block title}{fg=aneliseAzul,bg=aneliseGris!12!white}
\setbeamercolor{block body}{bg=aneliseGris!7!white}
\hypersetup{colorlinks=true,linkcolor=aneliseAzul,urlcolor=aneliseVerde}

% ---------- Comandos auxiliares ----------
\newcommand{\kurose}{Kurose \& Ross, \emph{Computer Networking: a Top-Down Approach}}
\newcommand{\forouzan}{Forouzan, \emph{Data Communications and Networking}}
\newcommand{\stallings}{Stallings, \emph{Wireless Communications \& Networks}}
\newcommand{\tanenbaum}{Tanenbaum, \emph{Computer Networks}}
\newcommand{\rfc}[1]{\href{https://www.rfc-editor.org/rfc/rfc#1.txt}{RFC~#1}}
% -----------------------------------------------------------------------------
%  Estilos TikZ comuns para desenhar redes (posicionamento relativo)
% -----------------------------------------------------------------------------
\tikzset{
  host/.style={draw,rounded corners=2mm,fill=aneliseAzul!15,inner sep=1.5mm,font=\scriptsize},
  rt/.style={draw,rounded corners=1mm,fill=aneliseLaranja!25,inner sep=1mm,font=\scriptsize},
  nube/.style={draw,cloud,aspect=2.2,fill=aneliseGris!15,inner sep=1.5mm,font=\scriptsize},
  el/.style={draw,rounded corners=1.2mm,fill=aneliseGris!18,inner sep=0.8mm,font=\scriptsize},
  enl/.style={thick},
  enlA/.style={thick,-Stealth},
} en cada aula.
%  Paleta de cores e estilos são definidos aqui uma única vez.
% =============================================================================

\usepackage[T1]{fontenc}
\usepackage{amsmath,amssymb,mathtools}
\usepackage{graphicx}
\usepackage{tikz}
\usetikzlibrary{positioning,arrows.meta,calc,backgrounds,decorations.pathmorphing,fit,shadows,shapes.symbols,shapes.geometric}
\usepackage{pgfplots}
\pgfplotsset{compat=1.16}
\usepackage{tabularx,booktabs,multirow,array}
\usepackage[normalem]{ulem}
\usepackage{hyperref}

% ---------- Paleta de cores própria ----------
\definecolor{aneliseAzul}{RGB}{20,60,110}
\definecolor{aneliseVerde}{RGB}{30,140,70}
\definecolor{aneliseLaranja}{RGB}{215,120,20}
\definecolor{aneliseGris}{RGB}{90,90,90}

% ---------- Tema ----------
\usetheme{Madrid}
\usecolortheme{whale}
\setbeamercolor{structure}{fg=aneliseAzul}
\setbeamercolor{title}{fg=aneliseAzul}
\setbeamercolor{frametitle}{fg=aneliseAzul}
\setbeamercolor{block title}{fg=aneliseAzul,bg=aneliseGris!12!white}
\setbeamercolor{block body}{bg=aneliseGris!7!white}
\hypersetup{colorlinks=true,linkcolor=aneliseAzul,urlcolor=aneliseVerde}

% ---------- Comandos auxiliares ----------
\newcommand{\kurose}{Kurose \& Ross, \emph{Computer Networking: a Top-Down Approach}}
\newcommand{\forouzan}{Forouzan, \emph{Data Communications and Networking}}
\newcommand{\stallings}{Stallings, \emph{Wireless Communications \& Networks}}
\newcommand{\tanenbaum}{Tanenbaum, \emph{Computer Networks}}
\newcommand{\rfc}[1]{\href{https://www.rfc-editor.org/rfc/rfc#1.txt}{RFC~#1}}
% -----------------------------------------------------------------------------
%  Estilos TikZ comuns para desenhar redes (posicionamento relativo)
% -----------------------------------------------------------------------------
\tikzset{
  host/.style={draw,rounded corners=2mm,fill=aneliseAzul!15,inner sep=1.5mm,font=\scriptsize},
  rt/.style={draw,rounded corners=1mm,fill=aneliseLaranja!25,inner sep=1mm,font=\scriptsize},
  nube/.style={draw,cloud,aspect=2.2,fill=aneliseGris!15,inner sep=1.5mm,font=\scriptsize},
  el/.style={draw,rounded corners=1.2mm,fill=aneliseGris!18,inner sep=0.8mm,font=\scriptsize},
  enl/.style={thick},
  enlA/.style={thick,-Stealth},
} en cada aula.
%  Paleta de cores e estilos são definidos aqui uma única vez.
% =============================================================================

\usepackage[T1]{fontenc}
\usepackage{amsmath,amssymb,mathtools}
\usepackage{graphicx}
\usepackage{tikz}
\usetikzlibrary{positioning,arrows.meta,calc,backgrounds,decorations.pathmorphing,fit,shadows,shapes.symbols,shapes.geometric}
\usepackage{pgfplots}
\pgfplotsset{compat=1.16}
\usepackage{tabularx,booktabs,multirow,array}
\usepackage[normalem]{ulem}
\usepackage{hyperref}

% ---------- Paleta de cores própria ----------
\definecolor{aneliseAzul}{RGB}{20,60,110}
\definecolor{aneliseVerde}{RGB}{30,140,70}
\definecolor{aneliseLaranja}{RGB}{215,120,20}
\definecolor{aneliseGris}{RGB}{90,90,90}

% ---------- Tema ----------
\usetheme{Madrid}
\usecolortheme{whale}
\setbeamercolor{structure}{fg=aneliseAzul}
\setbeamercolor{title}{fg=aneliseAzul}
\setbeamercolor{frametitle}{fg=aneliseAzul}
\setbeamercolor{block title}{fg=aneliseAzul,bg=aneliseGris!12!white}
\setbeamercolor{block body}{bg=aneliseGris!7!white}
\hypersetup{colorlinks=true,linkcolor=aneliseAzul,urlcolor=aneliseVerde}

% ---------- Comandos auxiliares ----------
\newcommand{\kurose}{Kurose \& Ross, \emph{Computer Networking: a Top-Down Approach}}
\newcommand{\forouzan}{Forouzan, \emph{Data Communications and Networking}}
\newcommand{\stallings}{Stallings, \emph{Wireless Communications \& Networks}}
\newcommand{\tanenbaum}{Tanenbaum, \emph{Computer Networks}}
\newcommand{\rfc}[1]{\href{https://www.rfc-editor.org/rfc/rfc#1.txt}{RFC~#1}}
% -----------------------------------------------------------------------------
%  Estilos TikZ comuns para desenhar redes (posicionamento relativo)
% -----------------------------------------------------------------------------
\tikzset{
  host/.style={draw,rounded corners=2mm,fill=aneliseAzul!15,inner sep=1.5mm,font=\scriptsize},
  rt/.style={draw,rounded corners=1mm,fill=aneliseLaranja!25,inner sep=1mm,font=\scriptsize},
  nube/.style={draw,cloud,aspect=2.2,fill=aneliseGris!15,inner sep=1.5mm,font=\scriptsize},
  el/.style={draw,rounded corners=1.2mm,fill=aneliseGris!18,inner sep=0.8mm,font=\scriptsize},
  enl/.style={thick},
  enlA/.style={thick,-Stealth},
}

\title[Redes sem Fio]{Redes sem Fio\\ 5G e Redes Móveis de Nova Geração}
\subtitle{Aula 7}
\author[Anelise Munaretto]{Profa. Anelise Munaretto}
\institute[Curso]{Redes sem Fio --- Ciência da Computação}
\date{2026}

\begin{document}

\begin{frame}[plain]
  \titlepage
\end{frame}

\begin{frame}{Agenda}
  \tableofcontents
\end{frame}

\begin{frame}{Objetivos de aprendizagem}
  Ao final desta aula, você será capaz de:
  \begin{itemize}
    \item \textbf{distinguir} as arquiteturas 5G SA e NSA e as implicações de cada uma na implantação;
    \item \textbf{explicar} o papel do NRF e das funções de rede na Service-Based Architecture (SBA);
    \item \textbf{calcular} a latência fim-a-fim de um enlace URLLC e a taxa espectral em mMIMO/mmWave;
    \item \textbf{comparar} os casos de uso eMBB, URLLC e mMTC e o papel do \emph{network slicing};
    \item \textbf{descrever} os caminhos de evolução rumo ao 6G (bandas THz, ISAC, redes 3D).
  \end{itemize}
  \begin{alertblock}{Pré-requisitos}
    \begin{itemize}
      \item Aula 2 (conceitos de enlace sem fio e QoS) e Aula 5 (5G NR, Wi-Fi e LPWAN).
    \end{itemize}
  \end{alertblock}
\end{frame}

\section{Evolução das redes móveis}

\begin{frame}{Evolução: de 1G a 5G/6G}
  \begin{center}
    \scriptsize
    \begin{tabular}{@{}llll@{}}
      \toprule
      \textbf{Geração} & \textbf{Período} & \textbf{Tecnologia} & \textbf{Taxa típica}\\
      \midrule
      1G & 1980s & AMPS, NMT (analógico) & ---\\
      2G & 1990s & GSM, CDMA (digital) & 64~kbps\\
      3G & 2000s & UMTS, CDMA2000 & 2~Mbps\\
      4G (LTE) & 2010s & OFDMA, MIMO & 150~Mbps\\
      5G NR & 2020s & OFDM flex., mMIMO & 1--20~Gbps\\
      5G-Advanced & 2024+ & IA/ML na rede & $>$ 10~Gbps\\
      6G (pesquisa) & 2030+ & THz, sensing integrado & $>$ 100~Gbps\\
      \bottomrule
    \end{tabular}
  \end{center}
  \begin{alertblock}{Salto 4G $\rightarrow$ 5G}
    Não apenas mais banda: latência de 1~ms, densidade de 10$^6$ dispositivos/km$^2$, \emph{network slicing}, \emph{edge computing}.
  \end{alertblock}
\end{frame}

\section{Espectro e numerologia}

\begin{frame}{Bandas 5G: FR1 (sub-6~GHz) e FR2 (mmWave)}
  \begin{center}
    \scriptsize
    \begin{tabular}{@{}p{1.5cm}p{3.3cm}p{7.4cm}@{}}
      \toprule
      \textbf{Faixa} & \textbf{Frequência} & \textbf{Bandas NR típicas (3GPP TS 38.101)}\\
      \midrule
      FR1 & 410~MHz--7,125~GHz & n41 (2,5~GHz), n77 (3,3--4,2~GHz), \textbf{n78} (3,3--3,8~GHz)\\
      FR2-1 & 24,25--52,6~GHz & \textbf{n257} (26,5--29,5~GHz), \textbf{n258} (24,25--27,5~GHz), \textbf{n260} (37--40~GHz), \textbf{n261} (27,5--28,35~GHz)\\
      FR2-2 & 52,6--71~GHz & n262 (47,2--48,2~GHz), n263 (57--71~GHz), bandas de 60~GHz (Rel-17)\\
      \bottomrule
    \end{tabular}
  \end{center}
  \begin{columns}
    \begin{column}{0.5\textwidth}
      \begin{block}{FR1 --- cobertura}
        \begin{itemize}
          \item propagação próxima do 4G, boa difração
          \item largura de banda por portadora até 100~MHz
          \item \emph{mid-band} (3,5~GHz) é o ``cavalo de batalha''
        \end{itemize}
      \end{block}
    \end{column}
    \begin{column}{0.5\textwidth}
      \begin{alertblock}{FR2 --- capacidade}
        \begin{itemize}
          \item largura de banda de 400~MHz a 2~GHz
          \item alta atenuação (O$_2$, chuva, bloqueio)
          \item \emph{beamforming} obrigatório; curto alcance
        \end{itemize}
      \end{alertblock}
    \end{column}
  \end{columns}
  \note{n78 (3,5 GHz) é a principal banda FR1 de implantação; FR2 usa n257/n258/n260/n261. FR2-2 (60 GHz) chega no Rel-17.}
\end{frame}

\begin{frame}{Numerologia OFDM e duração do slot}
  \begin{columns}
    \begin{column}{0.46\textwidth}
      \begin{center}
        \scriptsize
        \begin{tabular}{@{}cccc@{}}
          \toprule
          $\mu$ & SCS & Slots/subframe & Duração do slot\\
          \midrule
          0 & 15~kHz & 1 & 1~ms\\
          1 & 30~kHz & 2 & 0,5~ms\\
          2 & 60~kHz & 4 & 0,25~ms\\
          3 & 120~kHz & 8 & 125~$\mu$s\\
          4 & 240~kHz & 16 & 62,5~$\mu$s\\
          \bottomrule
        \end{tabular}
      \end{center}
      \[
        T_{\text{slot}} = \frac{1~\text{ms}}{2^{\mu}}, \qquad
        \text{SCS} = 15\cdot 2^{\mu}~\text{kHz}
      \]
    \end{column}
    \begin{column}{0.54\textwidth}
      \begin{center}
        \begin{tikzpicture}[font=\tiny,>=Stealth]
          \draw[fill=aneliseAzul!15,draw=aneliseAzul] (0,1.6) rectangle (10,2.2);
          \foreach \i in {1,...,9} \draw[aneliseAzul] (\i,1.6)--(\i,2.2);
          \node[right,font=\scriptsize] at (10.1,1.9) {radio frame = 10~ms};
          \draw[fill=aneliseVerde!18,draw=aneliseVerde] (0,0.7) rectangle (1,1.2);
          \node[right,font=\scriptsize] at (1.1,0.95) {1 subframe = 1~ms};
          \draw[fill=aneliseLaranja!20,draw=aneliseLaranja] (0,0.1) rectangle (0.5,0.5);
          \draw[fill=aneliseLaranja!20,draw=aneliseLaranja] (0.5,0.1) rectangle (1,0.5);
          \node[right,font=\scriptsize] at (1.1,0.3) {2$^\mu$ slots/subframe ($\mu=1$: 0,5~ms)};
        \end{tikzpicture}
      \end{center}
    \end{column}
  \end{columns}
  \begin{block}{Por que flexibilizar?}
    15~kHz para \textbf{cobertura} (eMBB rural, mMTC); 60--120~kHz para \textbf{latência baixa} (URLLC) e menor \emph{overhead} relativo de CP (4,7~$\mu$s a 15~kHz).
  \end{block}
  \note{O prefixo cíclico (CP) normal é de 4,7 $\mu$s a 15 kHz e escala inversamente com a SCS.}
\end{frame}

\begin{frame}{Duplexação e formatos de slot}
  \begin{columns}
    \begin{column}{0.5\textwidth}
      \begin{block}{FDD (paired spectrum)}
        \begin{itemize}
          \item UL e DL em frequências separadas
          \item \emph{full-duplex} por banda
          \item bom para cobertura e latência simétrica
          \item espectro ``casado'' mais caro
        \end{itemize}
      \end{block}
    \end{column}
    \begin{column}{0.5\textwidth}
      \begin{alertblock}{TDD (unpaired spectrum)}
        \begin{itemize}
          \item mesma banda, separação no \textbf{tempo}
          \item 3GPP TS 38.213: \emph{slot formats} D (downlink), U (uplink) e S (flexível)
          \item \emph{reciprocidade de canal} favorece mMIMO
          \item exige sincronização e \emph{guard periods}
        \end{itemize}
      \end{alertblock}
    \end{column}
  \end{columns}
  \begin{center}
    \begin{tikzpicture}[font=\tiny]
      \foreach \i/\t in {0/D,1/D,2/D,3/S,4/U,5/U,6/D,7/D,8/D,9/S} {
        \draw[fill=aneliseAzul!18,draw=aneliseAzul] (\i,0) rectangle ++(0.9,0.6);
        \node at (\i+0.45,0.3) {\t};
      }
      \node[right,font=\scriptsize] at (9.1,0.3) {TDD: padrão D/U/S (10~ms)};
    \end{tikzpicture}\\[1mm]
    {\scriptsize \textbf{Figura 1:} configuração TDD com slots de downlink (D), uplink (U) e flexível (S).}
  \end{center}
  \note{Em TDD, o padrão D/U/S é semi-estático (SIB1) e pode ser adaptado dinamicamente por DCI.}
\end{frame}

\section{Arquitetura 5G}

\begin{frame}{Arquitetura 5G: SA vs.\ NSA}
  \begin{columns}
    \begin{column}{0.5\textwidth}
      \textbf{NSA (Non-Standalone)}
      \begin{itemize}
        \item usa 5G NR apenas no \emph{data plane}
        \item \emph{control plane} ainda pelo 4G (eNB)
        \item implantação rápida; 3GPP Rel.~15
      \end{itemize}
      \textbf{SA (Standalone)}
      \begin{itemize}
        \item 5G NR tanto em controle quanto em dados
        \item new 5G core (SBA --- Service-Based Architecture)
        \item permite \emph{slicing}, URLLC, mMTC
      \end{itemize}
    \end{column}
    \begin{column}{0.5\textwidth}
      \begin{center}
        \begin{tikzpicture}[font=\scriptsize,
            nf/.style={draw,rounded corners=2mm,fill=aneliseAzul!13,inner sep=1mm,text width=2.6cm,align=center},
            nc/.style={draw,rounded corners=2mm,fill=aneliseLaranja!20,inner sep=1mm,text width=2.6cm,align=center},
          ]
          \node[nf](ue){UE};
          \node[nc,right=1.6cm of ue](gnb){gNB (5G)};
          \node[nf,right=1.6cm of gnb](cn){5GC (SBA)};
          \draw[<->,thick,aneliseVerde] (ue) -- (gnb);
          \draw[<->,thick,aneliseVerde] (gnb) -- (cn);
          \node[below=1mm of cn,font=\tiny,text width=3.2cm,align=center]{SA: controle e dados sobre 5G};
          \node[draw,fill=aneliseGris!15,rounded corners=1.5mm,below=1.6cm of gnb](leg){Arquitetura 5G SA simples};
        \end{tikzpicture}\\[1mm]
        {\scriptsize \textbf{Figura 2:} arquitetura 5G SA: UE conectado ao gNB (5G NR), que se liga ao núcleo 5GC baseado em SBA.}
      \end{center}
    \end{column}
  \end{columns}
  \note{Reforçar que no NSA o controle fica no eNB 4G e a implantação é mais rápida; o SA é o caminho para slicing e URLLC.}
\end{frame}

\begin{frame}{SBA --- Service-Based Architecture}
  \begin{itemize}
    \item Core 5G baseado em \textbf{serviços} (não em nós fixos)
    \item Funções de rede (NF) expõem APIs RESTful
      \begin{itemize}
        \item AMF (mobilidade), SMF (sessão), UPF (pacotes), NRF (registro), AUSF (autenticação), \ldots
      \end{itemize}
    \item \textbf{Interfaces abertas}: desacoplamento entre funções
    \item Facilita \emph{network slicing} e implantação em nuvem
  \end{itemize}
  \begin{center}
    \begin{tikzpicture}[font=\scriptsize,
        nf/.style={draw,rounded corners=1.5mm,fill=aneliseAzul!13,inner sep=0.8mm,minimum width=1.8cm,align=center},
      ]
      \node[nf](amf){AMF};
      \node[nf,right=1cm of amf](smf){SMF};
      \node[nf,right=1cm of smf](upf){UPF};
      \node[nf,above=1cm of amf](nrf){NRF};
      \node[nf,above=1cm of smf](ausf){AUSF};
      \node[nf,below=1cm of upf](pcf){PCF};
      \draw[thick,aneliseGris] (amf) -- (smf) -- (upf);
      \draw[thick,aneliseGris] (nrf) -- (amf);
      \draw[thick,aneliseGris] (ausf) -- (smf);
      \draw[thick,aneliseGris] (pcf) -- (upf);
      \node[draw,fill=aneliseVerde!12,rounded corners=2mm,below=4mm of pcf](bus){\texttt{bus de serviço}};
    \end{tikzpicture}\\[1mm]
    {\scriptsize \textbf{Figura 3:} funções de rede do 5GC (AMF, SMF, UPF, NRF, AUSF, PCF) interconectadas pelo \emph{bus} de serviço SBA.}
  \end{center}
  \note{O NRF é o ``catálogo'' de serviços: cada NF se registra e descobre as demais via APIs RESTful; isso desacopla as funções e viabiliza o slicing.}
\end{frame}

\section{Arquitetura avançada}

\begin{frame}{Dual connectivity: EN-DC, NR-DC e NE-DC}
  \begin{columns}
    \begin{column}{0.5\textwidth}
      \begin{itemize}
        \item UE mantém \textbf{duas conexões} simultâneas: nó mestre (MN) e nó secundário (SN)
        \item agregação de portadoras entre tecnologias/nós
        \item \textbf{EN-DC} (Option 3): MN = eNB (LTE), SN = gNB (NR) --- base do NSA
        \item \textbf{NR-DC}: MN e SN são gNB (SA)
        \item \textbf{NE-DC}: MN = gNB (NR), SN = eNB (LTE)
      \end{itemize}
    \end{column}
    \begin{column}{0.5\textwidth}
      \begin{center}
        \begin{tikzpicture}[font=\scriptsize]
          \node[draw,rounded corners,fill=aneliseAzul!15,minimum width=1.2cm](ue){UE};
          \node[draw,rounded corners,fill=aneliseLaranja!20,right=2.2cm of ue](mn){MN (eNB/gNB)};
          \node[draw,rounded corners,fill=aneliseVerde!18,above=0.9cm of mn](sn){SN (gNB)};
          \draw[<->,thick,aneliseAzul] (ue) -- (mn);
          \draw[<->,thick,aneliseVerde] (ue) -- (sn);
          \draw[<->,thick,aneliseGris] (mn) -- (sn);
          \node[font=\tiny,below=0.15cm of mn]{X2/Xn};
        \end{tikzpicture}\\[1mm]
        {\scriptsize \textbf{Figura 4:} dual connectivity: o UE agrega o nó mestre e o secundário.}
      \end{center}
    \end{column}
  \end{columns}
  \begin{block}{Benefícios}
    Pico de taxa (agregação), ancoragem de mobilidade no MN e implantação incremental do NR sobre o LTE (NSA).
  \end{block}
\end{frame}

\begin{frame}{PDU sessions, QoS flows e DRBs}
  \begin{itemize}
    \item \textbf{PDU Session}: associação UE $\leftrightarrow$ rede de dados (DN), com IP (ou Ethernet/não-IP); ancorada no UPF
    \item cada PDU session contém um ou mais \textbf{QoS Flows} (identificados pelo \textbf{QFI})
    \item \textbf{5QI}: índice de QoS (tipo de recurso GBR/non-GBR, prioridade, \emph{packet delay budget}, taxa de erro)
    \item QoS Flow $\rightarrow$ \textbf{DRB} (Data Radio Bearer) no rádio, com mapeamento pelo gNB
  \end{itemize}
  \begin{center}
    \begin{tikzpicture}[font=\tiny,>=Stealth]
      \node[draw,rounded corners,fill=aneliseAzul!15,minimum width=1.1cm](ue){UE};
      \node[draw,rounded corners,fill=aneliseLaranja!20,right=2.0cm of ue,minimum width=1.3cm](gnb){gNB};
      \node[draw,rounded corners,fill=aneliseVerde!18,right=2.0cm of gnb,minimum width=1.3cm](upf){UPF};
      \node[draw,rounded corners,fill=aneliseGris!18,right=1.4cm of upf](dn){DN};
      \draw[thick,aneliseLaranja] ($(ue.east)+(0,0.5)$) -- ($(gnb.west)+(0,0.5)$) -- ($(upf.west)+(0,0.5)$) -- ($(dn.west)+(0,0.5)$);
      \draw[thick,aneliseVerde] ($(ue.east)+(0,0)$) -- ($(gnb.west)+(0,0)$) -- ($(upf.west)+(0,0)$) -- ($(dn.west)+(0,0)$);
      \draw[thick,aneliseAzul] ($(ue.east)+(0,-0.5)$) -- ($(gnb.west)+(0,-0.5)$) -- ($(upf.west)+(0,-0.5)$) -- ($(dn.west)+(0,-0.5)$);
      \node[above,font=\tiny,text=aneliseLaranja] at ($(gnb.west)!0.5!(upf.west)+(0,0.5)$){5QI 1 (GBR, voz)};
      \node[font=\tiny,text=aneliseVerde] at ($(gnb.west)!0.5!(upf.west)+(0,0)$){5QI 9 (non-GBR, eMBB)};
      \node[below,font=\tiny,text=aneliseAzul] at ($(gnb.west)!0.5!(upf.west)+(0,-0.5)$){5QI 82 (URLLC)};
      \node[below,font=\tiny] at ($(ue.east)!0.5!(gnb.west)+(0,-0.9)$){DRBs};
      \node[below,font=\tiny] at ($(gnb.east)!0.5!(upf.west)+(0,-0.9)$){N3/GTP-U};
    \end{tikzpicture}\\[1mm]
    {\scriptsize \textbf{Figura 5:} uma PDU session transporta três QoS flows (voz GBR, eMBB e URLLC) sobre DRBs distintos.}
  \end{center}
  \note{N3 = interface gNB--UPF (GTP-U); N4 = SMF--UPF; N6 = UPF--DN. O mapeamento QoS flow $\to$ DRB é decidido pelo gNB.}
\end{frame}

\begin{frame}{Mobilidade e handover no 5G}
  \begin{columns}
    \begin{column}{0.52\textwidth}
      \begin{itemize}
        \item \textbf{Intra-gNB}: troca de feixe/célula sem mudar o nó
        \item \textbf{Xn-based} (inter-gNB): sinalização direta gNB--gNB (interface Xn); \emph{path switch} no UPF
        \item \textbf{N2-based} (inter-AMF): quando não há Xn; sinalização via núcleo (AMF)
        \item \textbf{CP/UP split}: controle no RAN/AMF, dados no UPF; a sessão PDU pode ser realocada
        \item \textbf{DAPS} (Rel-16): \emph{dual active protocol stack} para interrupção $\approx 0$~ms
      \end{itemize}
    \end{column}
    \begin{column}{0.48\textwidth}
      \begin{center}
        \begin{tikzpicture}[font=\tiny,>=Stealth]
          \node[draw,rounded corners,fill=aneliseAzul!15](src){gNB fonte};
          \node[draw,rounded corners,fill=aneliseVerde!18,right=1.4cm of src](tgt){gNB alvo};
          \node[draw,rounded corners,fill=aneliseLaranja!20] at ($(src)!0.5!(tgt)+(0,-1.4)$)(amf){AMF/UPF};
          \draw[<->,thick,aneliseGris] (src) -- node[above,font=\tiny]{Xn} (tgt);
          \draw[<->,thick] (src) -- (amf);
          \draw[<->,thick] (tgt) -- (amf);
          \node[font=\tiny,right=0.1cm of amf]{N2/N3};
        \end{tikzpicture}\\[1mm]
        {\scriptsize \textbf{Figura 6:} handover Xn (direto) e N2 (via núcleo).}
      \end{center}
    \end{column}
  \end{columns}
\end{frame}

\begin{frame}{MEC: arquitetura e casos de uso}
  \begin{columns}
    \begin{column}{0.5\textwidth}
      \begin{itemize}
        \item \textbf{ETSI MEC}: \emph{host} MEC na borda, orquestrador e \emph{MEC apps}
        \item \textbf{Local breakout}: UPF local (ULCL) desvia o tráfego para a borda
        \item latência típica $<$ 10~ms; economia de \emph{backhaul}
      \end{itemize}
      \begin{block}{Casos de uso}
        AR/VR, V2X, \emph{cloud gaming}, análise de vídeo, automação industrial, \emph{caching} de conteúdo.
      \end{block}
    \end{column}
    \begin{column}{0.5\textwidth}
      \begin{center}
        \begin{tikzpicture}[font=\tiny,>=Stealth]
          \node[draw,rounded corners,fill=aneliseAzul!15](ue){UE};
          \node[draw,rounded corners,fill=aneliseLaranja!20,right=1.4cm of ue](gnb){gNB};
          \node[draw,rounded corners,fill=aneliseVerde!18,above=0.9cm of gnb](mec){MEC host};
          \node[draw,rounded corners,fill=aneliseGris!18,right=1.6cm of gnb](upf){UPF};
          \node[draw,cloud,cloud puffs=9,aspect=2,fill=aneliseGris!12,right=1.2cm of upf,font=\tiny](cl){nuvem};
          \draw[<->,thick] (ue) -- (gnb);
          \draw[<->,thick,aneliseVerde] (gnb) -- (mec);
          \draw[<->,thick] (gnb) -- (upf);
          \draw[<->,thick] (upf) -- (cl);
        \end{tikzpicture}\\[1mm]
        {\scriptsize \textbf{Figura 7:} MEC na borda da RAN, com \emph{local breakout} no UPF.}
      \end{center}
    \end{column}
  \end{columns}
\end{frame}

\section{Tecnologias-chave}

\begin{frame}{Tecnologias-chave do 5G NR}
  \begin{columns}
    \begin{column}{0.48\textwidth}
      \textbf{mmWave (ondas milimétricas)}
      \begin{itemize}
        \item bandas 24--28~GHz, 39~GHz, 60~GHz
        \item grande largura de banda ($\times$10) mas alta atenuação
        \item \emph{beamforming} e \emph{beam management} obrigatórios
      \end{itemize}
      \textbf{Massive MIMO}
      \begin{itemize}
        \item dezenas/milhares de antenas (64T64R, 128T128R)
        \item multiplexação espacial (MU-MIMO)
        \item ganho de array, precisão de ângulo
      \end{itemize}
    \end{column}
    \begin{column}{0.48\textwidth}
      \textbf{OFDM flexível (numerologia)}
      \begin{itemize}
        \item subportadoras de 15/30/60/120~kHz
        \item \emph{slot} de 0,5~ms a 62,5~$\mu$s
      \end{itemize}
      \textbf{Network Slicing}
      \begin{itemize}
        \item redes virtuais sobre a mesma infraestrutura
        \item eMBB (banda larga), URLLC (latência), mMTC (IoT)
        \item isolamento de recursos por \emph{slice}
      \end{itemize}
      \textbf{MEC --- Multi-access Edge Computing}
      \begin{itemize}
        \item computação na borda da rede RAN
        \item latência $<$ 10~ms para aplicações críticas
      \end{itemize}
    \end{column}
  \end{columns}
\end{frame}

\begin{frame}{Casos de uso: eMBB, URLLC, mMTC}
  \begin{center}
    \begin{tikzpicture}[font=\scriptsize,
        c/.style={draw,rounded corners=2mm,fill=aneliseAzul!13,text width=3.4cm,align=center,inner sep=1.2mm},
      ]
      \node[c](embb){\textbf{eMBB}\\ Enhanced Mobile Broadband\\ $>$ 1~Gbps\\ streaming 8K, VR/AR};
      \node[c,right=2.6cm of embb](urllc){\textbf{URLLC}\\ Ultra-Reliable Low Latency\\ 1~ms, 99,999\%\\ indústria 4.0, veículos autônomos};
      \node[c,right=2.6cm of urllc](mmtc){\textbf{mMTC}\\ Massive Machine-Type\\ 10$^6$ disp./km$^2$\\ IoT massivo, sensores};
      \draw[thick,aneliseLaranja] (embb.east) -- (urllc.west);
      \draw[thick,aneliseGris] (urllc.east) -- (mmtc.west);
    \end{tikzpicture}\\[1mm]
    {\scriptsize \textbf{Figura 8:} os três casos de uso do 5G definidos pelo IMT-2020: eMBB, URLLC e mMTC.}
  \end{center}
  \begin{block}{5G-Advanced (3GPP Rel.\ 18, 2024)}
    \begin{itemize}
      \item IA/ML na otimização da rede (RM-NET)
      \item \emph{RedCap}: IoT de média taxa (Reduced Capability)
      \item NR-U (5G em espectro não licenciado)
      \item NTN (satélite) integrado ao 5G
    \end{itemize}
  \end{block}
\end{frame}

\section{mMIMO e beamforming}

\begin{frame}{Massive MIMO: princípios e ganhos}
  \begin{columns}
    \begin{column}{0.5\textwidth}
      \begin{itemize}
        \item dezenas a centenas de elementos de antena no gNB (64T64R, 128T128R, 256+)
        \item \textbf{MU-MIMO}: vários UEs no mesmo recurso tempo-frequência via \emph{precoding}
        \item \textbf{ganho de array}: $G = 10\log_{10}(N)$
        \item \emph{channel hardening} e \emph{favorable propagation}
      \end{itemize}
    \end{column}
    \begin{column}{0.5\textwidth}
      \begin{block}{Exemplos de ganho de array}
        \begin{center}
          \begin{tabular}{@{}cc@{}}
            \toprule
            $N$ antenas & $10\log_{10}N$\\
            \midrule
            16 & 12,0~dB\\
            64 & 18,1~dB\\
            256 & 24,1~dB\\
            1024 & 30,1~dB\\
            \bottomrule
          \end{tabular}
        \end{center}
      \end{block}
      \begin{exampleblock}{Cálculo}
        $N=64$: $10\log_{10}(64)=10\times 1{,}806=18{,}06~\text{dB}\approx \mathbf{18~dB}$.
      \end{exampleblock}
    \end{column}
  \end{columns}
  \note{Os 18 dB do exemplo compensam parte da atenuação adicional em mmWave; esse ganho é uma das razões do mMIMO em FR2.}
\end{frame}

\begin{frame}{Beamforming: analógico, digital e híbrido}
  \begin{center}
    \begin{tikzpicture}[font=\tiny,>=Stealth]
      \foreach \y in {0,0.22,...,1.54} \draw[fill=aneliseAzul!40,draw=aneliseAzul] (0,\y) rectangle ++(0.14,0.14);
      \node[below,font=\scriptsize] at (0.6,-0.2){Analógico};
      \draw[aneliseLaranja,thick] (0.2,0.85) -- (2.0,1.5);
      \draw[aneliseLaranja,thick] (0.2,0.85) -- (2.0,0.85);
      \draw[aneliseLaranja,thick] (0.2,0.85) -- (2.0,0.2);
      \node[right,font=\tiny] at (2.1,0.85){1 cadeia RF};
      \begin{scope}[xshift=4.5cm]
        \foreach \y in {0,0.22,...,1.54} \draw[fill=aneliseVerde!40,draw=aneliseVerde] (0,\y) rectangle ++(0.14,0.14);
        \node[below,font=\scriptsize] at (0.6,-0.2){Digital};
        \draw[aneliseVerde,thick] (0.2,1.0) -- (2.0,1.6);
        \draw[aneliseVerde,thick] (0.2,1.0) -- (2.0,1.2);
        \draw[aneliseVerde,thick] (0.2,1.0) -- (2.0,0.6);
        \draw[aneliseVerde,thick] (0.2,1.0) -- (2.0,0.3);
        \node[right,font=\tiny] at (2.1,1.0){1 RF/antena};
      \end{scope}
      \begin{scope}[xshift=9.0cm]
        \foreach \x in {0,0.22,...,0.66} \foreach \y in {0,0.22,...,1.54} \draw[fill=aneliseLaranja!40,draw=aneliseLaranja] (\x,\y) rectangle ++(0.14,0.14);
        \node[below,font=\scriptsize] at (0.6,-0.2){Híbrido};
        \draw[aneliseLaranja,thick] (0.2,1.0) -- (2.0,1.5);
        \draw[aneliseLaranja,thick] (0.2,1.0) -- (2.0,0.5);
        \node[right,font=\tiny] at (2.1,1.0){sub-arrays};
      \end{scope}
    \end{tikzpicture}\\[1mm]
    {\scriptsize \textbf{Figura 9:} comparativo de arquiteturas de \emph{beamforming} em mmWave.}
  \end{center}
  \begin{columns}
    \begin{column}{0.33\textwidth}
      \begin{block}{Analógico}
        Um feixe por vez, deslocadores de fase; barato e comum em FR2.
      \end{block}
    \end{column}
    \begin{column}{0.33\textwidth}
      \begin{block}{Digital}
        Feixes múltiplos e flexíveis; alto custo/consumo (1 RF por antena).
      \end{block}
    \end{column}
    \begin{column}{0.33\textwidth}
      \begin{block}{Híbrido}
        Compromisso: poucas cadeias RF com sub-arrays; padrão prático do 5G mmWave.
      \end{block}
    \end{column}
  \end{columns}
\end{frame}

\begin{frame}{Beam management}
  \begin{columns}
    \begin{column}{0.55\textwidth}
      \begin{itemize}
        \item \textbf{SSB} (\emph{Synchronization Signal Block}): varredura de feixes de sincronização (\emph{beam sweeping})
        \item \textbf{CSI-RS}: medição de canal para seleção/refinamento de feixe
        \item procedimentos P1/P2/P3: seleção de feixe do gNB e do UE
        \item \textbf{Beam failure recovery (BFR)}: detecção de falha e recuperação via RACH
        \item reporte (\emph{beam reporting}) por L1/L2
      \end{itemize}
    \end{column}
    \begin{column}{0.45\textwidth}
      \begin{center}
        \begin{tikzpicture}[font=\tiny,>=Stealth]
          \node[draw,rounded corners,fill=aneliseAzul!15](gnb){gNB};
          \draw[aneliseLaranja,thick] (gnb.east) -- ++(0.6,0.9);
          \draw[aneliseLaranja,thick] (gnb.east) -- ++(0.8,0.3);
          \draw[aneliseLaranja,thick] (gnb.east) -- ++(0.8,-0.3);
          \draw[aneliseLaranja,thick] (gnb.east) -- ++(0.6,-0.9);
          \node[draw,rounded corners,fill=aneliseVerde!18,right=2.2cm of gnb,yshift=0.3cm](ue){UE};
          \draw[<->,thick,aneliseGris] (gnb.east) -- (ue.west);
          \node[font=\tiny,below=0.1cm of gnb]{beam sweeping};
        \end{tikzpicture}\\[1mm]
        {\scriptsize \textbf{Figura 10:} varredura de feixes (SSB) e alinhamento com o UE.}
      \end{center}
    \end{column}
  \end{columns}
  \note{SSB = PSS + SSS + PBCH; até 64 feixes de SSB em FR2. O BFR usa o acesso aleatório para recuperar o feixe.}
\end{frame}

\section{5G e sem fio}

\begin{frame}{5G NR e Wi-Fi: convergência}
  \begin{itemize}
    \item 5G NR-U compete e coopera com Wi-Fi 6/6E na banda de 6~GHz
    \item \textbf{ATSSS} (Access Traffic Steering, Switching, Splitting): tráfego pode ser dividido entre 5G e Wi-Fi simultaneamente
    \item \textbf{5G LAN}: redes privativas 5G para indústria (substituindo Wi-Fi/Ethernet)
    \item \textbf{Comparação prática}:
      \begin{itemize}
        \item Wi-Fi 7: mais barato, ideal para ambientes internos densos
        \item 5G NR: mobilidade, \emph{slicing}, QoS fim-a-fim, cobertura externa
      \end{itemize}
  \end{itemize}
\end{frame}

\section{Espectro não licenciado e IoT}

\begin{frame}{NR-U e LAA: 5G em espectro não licenciado}
  \begin{itemize}
    \item \textbf{LAA} (LTE Licensed-Assisted Access): LTE em 5~GHz assistido por portadora licenciada
    \item \textbf{NR-U} (Rel-16): NR em 5~GHz e 6~GHz, modos \emph{standalone} e assistido
    \item \textbf{LBT} (\emph{Listen Before Talk}): obrigatório por regulação (ETSI/FCC) para coexistência justa
    \item coexistência com Wi-Fi 6/6E/7 na mesma banda; largura de canal de 20 a 80/100~MHz
  \end{itemize}
  \begin{columns}
    \begin{column}{0.5\textwidth}
      \begin{block}{Vantagens}
        Mais espectro disponível, \emph{offload} de tráfego e \emph{carrier aggregation}.
      \end{block}
    \end{column}
    \begin{column}{0.5\textwidth}
      \begin{alertblock}{Desafios}
        LBT introduz latência variável; sem garantias rígidas de QoS; interferência com WLAN.
      \end{alertblock}
    \end{column}
  \end{columns}
  \note{A banda de 6 GHz (5925--7125 MHz) é chave para o NR-U e para o Wi-Fi 6E/7.}
\end{frame}

\begin{frame}{RedCap: IoT de média capacidade (Rel-17)}
  \begin{columns}
    \begin{column}{0.5\textwidth}
      \begin{itemize}
        \item \emph{Reduced Capability} NR: simplifica o modem para reduzir custo e consumo
        \item 1 ou 2 receptores; largura de banda reduzida (20~MHz em FR1)
        \item \emph{half-duplex} FDD; menor pico e menor complexidade
        \item preenche a lacuna entre eMBB e mMTC
      \end{itemize}
    \end{column}
    \begin{column}{0.5\textwidth}
      \begin{block}{Casos de uso}
        \emph{Wearables}, sensores industriais, câmeras de vídeo e \emph{smart grid}.
      \end{block}
      \begin{exampleblock}{Números (Rel-17)}
        Pico de $\approx$ 100~Mbps DL / 50~Mbps UL; latência de dezenas de ms.
      \end{exampleblock}
    \end{column}
  \end{columns}
  \note{O RedCap é especificado no 3GPP TS 38.306 (Rel-17).}
\end{frame}

\begin{frame}{Integração NTN (Rel-17/18)}
  \begin{itemize}
    \item \textbf{NTN --- Non-Terrestrial Networks}: satélites LEO/GEO, HAPS e \emph{drones}
    \item Rel-17: primeiras especificações de \emph{NR NTN} e \emph{IoT NTN} (3GPP TR 38.811)
    \item Rel-18: \emph{enhancements} --- cobertura, mobilidade e \emph{timing advance} para LEO
    \item desafios: atraso de propagação longo, \emph{Doppler} alto e células em movimento
  \end{itemize}
  \begin{exampleblock}{Cobertura 3D}
    NTN estende o 5G a áreas remotas, marítimas e aéreas; é a base da futura rede 3D do 6G.
  \end{exampleblock}
  \note{Detalhes de NTN serão aprofundados na Aula 10.}
\end{frame}

\section{5G-Advanced}

\begin{frame}{5G-Advanced (Rel-18/19): rumo ao 6G}
  \begin{columns}
    \begin{column}{0.5\textwidth}
      \begin{block}{Rel-18 (2024)}
        \begin{itemize}
          \item \textbf{IA/ML na RAN}: feedback de CSI, \emph{beam management} e posicionamento
          \item \textbf{eficiência energética} de rede (\emph{network energy savings})
          \item \textbf{MBS}: multicast/broadcast de serviços
          \item RedCap e NR-U evoluídos; NTN aprimorado
        \end{itemize}
      \end{block}
    \end{column}
    \begin{column}{0.5\textwidth}
      \begin{alertblock}{Rel-19 e além}
        \begin{itemize}
          \item IA/ML nativo e \emph{air interface} com \emph{sensing}
          \item \emph{Ambient IoT} e posicionamento de alta precisão
          \item convergência com computação de borda e nuvem
        \end{itemize}
      \end{alertblock}
    \end{column}
  \end{columns}
  \begin{center}
    {\scriptsize Rel-18 é o primeiro lançamento rotulado \textbf{5G-Advanced}; o Rel-21 (2028) iniciará o IMT-2030 (6G).}
  \end{center}
\end{frame}

\section{Rumo ao 6G}

\begin{frame}{Rumo ao 6G (2030+)}
  \begin{columns}
    \begin{column}{0.5\textwidth}
      \textbf{Bandas THz (0,1--10 THz)}
      \begin{itemize}
        \item largura de banda extrema ($>$ 100~Gbps)
        \item comunicação + \emph{sensing} integrados (ISAC)
        \item desafios: alcance, hardware, regulação
      \end{itemize}
      \textbf{IA nativa na rede}
      \begin{itemize}
        \item planejamento, alocação, predição com ML
        \item \emph{zero-touch} operações (ZSM)
      \end{itemize}
    \end{column}
    \begin{column}{0.5\textwidth}
      \textbf{Rede 3D integrada}
      \begin{itemize}
        \item terrestre + satélite LEO + \emph{drone base stations}
        \item cobertura global contínua
      \end{itemize}
      \textbf{Sustentabilidade}
      \begin{itemize}
        \item redes \emph{green}: eficiência energética por bit
        \item energy harvesting para IoT
      \end{itemize}
    \end{column}
  \end{columns}
  \begin{exampleblock}{Pesquisa aberta}
    3GPP Rel.\ 21 (2028) definirá os requisitos; ITU-R WP5D prepara o IMT-2030.
  \end{exampleblock}
  \note{Em 6G, o destaque é a convergência comunicação--sensoriamento (ISAC) e a integração com LEO/drones; o IMT-2030 será definido no 3GPP Rel.\ 21.}
\end{frame}

\begin{frame}{6G: RIS --- superfícies refletoras inteligentes}
  \begin{columns}
    \begin{column}{0.52\textwidth}
      \begin{itemize}
        \item \textbf{RIS} (\emph{Reconfigurable Intelligent Surface}): arranjo de elementos refletores passivos com fase controlável
        \item cria um \textbf{caminho virtual de LOS} ao contornar obstáculos
        \item não amplifica o sinal: apenas o redireciona (baixo consumo)
        \item aplicações: cobertura em mmWave/THz, \emph{dead zones}, eficiência energética
      \end{itemize}
    \end{column}
    \begin{column}{0.48\textwidth}
      \begin{center}
        \begin{tikzpicture}[font=\tiny,>=Stealth]
          \node[draw,rounded corners,fill=aneliseAzul!15](bs){gNB};
          \node[fill=aneliseGris!40,minimum width=0.22cm,minimum height=2.0cm] at (3,0.5){};
          \node[font=\tiny,below] at (3,0.0){obstáculo};
          \node[draw,fill=aneliseLaranja!30,minimum width=0.2cm,minimum height=1.4cm] at (5.2,1.8){};
          \node[font=\tiny,above] at (5.2,2.5){RIS};
          \node[draw,rounded corners,fill=aneliseVerde!18] at (6.8,0.2)(ue){UE};
          \draw[-Stealth,thick,aneliseAzul] (bs.east) -- (5.2,2.4);
          \draw[-Stealth,thick,aneliseVerde] (5.2,1.2) -- (ue.north);
          \draw[dashed,aneliseGris] (bs.east) -- (3,0.3);
        \end{tikzpicture}\\[1mm]
        {\scriptsize \textbf{Figura 11:} a RIS reflete o feixe do gNB ao redor do obstáculo até o UE.}
      \end{center}
    \end{column}
  \end{columns}
\end{frame}

\begin{frame}{6G: ISAC e a convergência comunicação--sensoriamento}
  \begin{itemize}
    \item \textbf{ISAC} (\emph{Integrated Sensing and Communication}): o mesmo sinal serve para comunicar e ``enxergar'' o ambiente
    \item o gNB funciona como radar: estima distância, velocidade e ângulo de alvos
    \item casos de uso: localização de alta precisão, mapeamento do ambiente, detecção de obstáculos, gestos e V2X
    \item \emph{trade-off}: dividir recursos entre taxa de comunicação e resolução de sensoriamento
  \end{itemize}
  \begin{columns}
    \begin{column}{0.5\textwidth}
      \begin{block}{Habilitadores}
        Bandas altas (mmWave/THz), grandes arrays (mMIMO) e IA para processamento do eco.
      \end{block}
    \end{column}
    \begin{column}{0.5\textwidth}
      \begin{alertblock}{Visão IMT-2030}
        Sensoriamento integrado é um dos novos casos de uso do 6G no \emph{framework} ITU-R M.2160.
      \end{alertblock}
    \end{column}
  \end{columns}
\end{frame}

\section{Implantação}

\begin{frame}{Densificação: macro células e small cells (HetNet)}
  \begin{columns}
    \begin{column}{0.5\textwidth}
      \begin{itemize}
        \item \textbf{macro cells}: cobertura ampla (km), alta potência, mobilidade
        \item \textbf{small cells} (micro/pico/femto): baixa potência, curto alcance, alta densidade
        \item \textbf{HetNet}: camadas sobrepostas para aumentar capacidade e \emph{area spectral efficiency}
        \item \emph{backhaul} das small cells: fibra, mmWave ou PLC
      \end{itemize}
    \end{column}
    \begin{column}{0.5\textwidth}
      \begin{center}
        \begin{tikzpicture}[font=\tiny]
          \draw[fill=aneliseAzul!10,draw=aneliseAzul] (0,0) circle (2.2);
          \node[font=\scriptsize,text=aneliseAzul] at (0,1.9){macro};
          \foreach \x/\y in {-1.1/0.5, 1.0/-0.8, 0.2/1.0} {
            \draw[fill=aneliseLaranja!25,draw=aneliseLaranja] (\x,\y) circle (0.55);
          }
          \foreach \x/\y in {-1.1/0.5, 1.0/-0.8, 0.2/1.0} {
            \node[circle,fill=aneliseLaranja,inner sep=1pt] at (\x,\y){};
          }
          \node[font=\tiny,text=aneliseLaranja] at (-1.1,1.35){small};
        \end{tikzpicture}\\[1mm]
        {\scriptsize \textbf{Figura 12:} HetNet: macro célula com small cells sobrepostas.}
      \end{center}
    \end{column}
  \end{columns}
  \note{A densificação é o principal vetor de capacidade do 5G; a interferência entre camadas exige coordenação (eICIC/ABS).}
\end{frame}

\begin{frame}{Desafios de implantação}
  \begin{columns}
    \begin{column}{0.5\textwidth}
      \begin{block}{\emph{Backhaul}}
        Conectar milhares de small cells; fibra nem sempre viável; alternativas em mmWave e malha.
      \end{block}
      \begin{block}{Energia}
        RAN consome a maior parte da energia da rede; modos de \emph{sleep} e fontes renováveis.
      \end{block}
    \end{column}
    \begin{column}{0.5\textwidth}
      \begin{alertblock}{Espectro}
        Licenciamento caro; \emph{mid-band} disputada; uso compartilhado (CBRS, 6~GHz).
      \end{alertblock}
      \begin{alertblock}{Cobertura indoor}
        mmWave não penetra bem; \emph{small cells} indoor, DAS e Wi-Fi como complemento.
      \end{alertblock}
    \end{column}
  \end{columns}
  \begin{center}
    {\scriptsize CAPEX/OPEX, aquisição de sites e regulação também condicionam o ritmo de implantação.}
  \end{center}
\end{frame}

\begin{frame}{5G NR vs.\ Wi-Fi 7: quando usar cada um}
  \begin{center}
    \scriptsize
    \begin{tabular}{@{}p{2.6cm}p{5.0cm}p{5.0cm}@{}}
      \toprule
      \textbf{Critério} & \textbf{5G NR} & \textbf{Wi-Fi 7 (IEEE 802.11be)}\\
      \midrule
      Mobilidade & \emph{handover} entre células, alta velocidade & limitada ao ESS; sem \emph{handover} celular\\
      QoS fim-a-fim & 5QI, \emph{QoS flows}, garantias e URLLC & EDCA/WMM, melhor esforço com prioridades\\
      \emph{Slicing} & sim, isolamento de recursos & não nativo (SSID/VLAN)\\
      Cobertura & macro (km) e indoor & dezenas de metros, ambiente controlado\\
      Espectro & licenciado (+ não licenciado NR-U) & não licenciado (2,4/5/6~GHz)\\
      Custo & CAPEX/OPEX altos; \emph{chipsets} caros & baixo; amplo ecossistema\\
      Casos de uso & eMBB/URLLC/mMTC, V2X, indústria & WLAN densa, casa, escritório, \emph{offload}\\
      \bottomrule
    \end{tabular}
  \end{center}
  \note{Wi-Fi 7 traz MLO, canais de 320 MHz e 4096-QAM; o 5G se destaca em mobilidade, QoS e slicing.}
\end{frame}

\begin{frame}{Exercício resolvido: latência fim-a-fim em URLLC}
  \textbf{Problema:} um quadro de \textbf{1000 bytes} é enviado por um enlace 5G NR a \textbf{1~Gbps}. O enlace tem \textbf{10~km} de comprimento. Estime a latência fim-a-fim (desprezando filas).
  \begin{block}{Resolução}
    \begin{enumerate}
      \item Tempo de transmissão (L = tamanho, R = taxa):
        \[
          T_{\text{tx}} = \frac{L}{R} = \frac{1000 \times 8~\text{bits}}{10^9~\text{bps}} = \frac{8000}{10^9} = 8 \times 10^{-6}~\text{s} = \mathbf{8~\mu s}.
        \]
      \item Tempo de propagação (d = 10 km, v $\approx c = 3\times10^8$ m/s):
        \[
          T_{\text{prop}} = \frac{d}{v} = \frac{10 \times 10^3~\text{m}}{3 \times 10^8~\text{m/s}} \approx 3,33 \times 10^{-5}~\text{s} \approx \mathbf{33,3~\mu s}.
        \]
      \item Latência fim-a-fim:
        \[
          T_{\text{e2e}} = T_{\text{tx}} + T_{\text{prop}} \approx 8 + 33,3 = \mathbf{41,3~\mu s}.
        \]
    \end{enumerate}
  \end{block}
  \begin{alertblock}{Leitura do resultado}
    O alvo URLLC é $\le 1$~ms fim-a-fim; só o enlace consome $\approx 41~\mu$s, sobrando orçamento para \emph{scheduling}, filas e processamento na borda (MEC).
  \end{alertblock}
\end{frame}

\begin{frame}{Resumo da aula}
  \begin{itemize}
    \item \textbf{SA vs NSA}: SA usa 5G NR no controle e nos dados sobre o 5GC (SBA); NSA mantém o controle no núcleo 4G.
    \item \textbf{SBA}: AMF, SMF, UPF, NRF, AUSF e PCF expõem serviços RESTful num barramento de serviço; o NRF registra e descobre as funções.
    \item \textbf{Tecnologias-chave}: mmWave (24--60~GHz) com \emph{beamforming}, mMIMO (64T64R+), numerologia OFDM flexível, \emph{network slicing} e MEC.
    \item \textbf{Casos de uso}: eMBB, URLLC e mMTC; 5G-Advanced (Rel.~18) adiciona IA/ML, RedCap, NR-U e NTN.
    \item \textbf{6G}: bandas THz, ISAC, redes 3D (LEO/drones) e IA nativa.
  \end{itemize}
\end{frame}

\begin{frame}{Autoavaliação}
  \begin{enumerate}
    \item Na arquitetura 5G SA, o plano de controle:
      \begin{itemize}
        \item (a) ainda é executado pelo eNB 4G (LTE)
        \item (b) usa exclusivamente a função UPF
        \item (c) é executado pelo 5G NR sobre o 5GC (SBA)
      \end{itemize}
      \pause
      \textbf{Gabarito:} (c)
    \item Na SBA, qual função de rede é responsável por registrar e descobrir os serviços das demais funções?
      \begin{itemize}
        \item (a) NRF
        \item (b) SMF
        \item (c) UPF
      \end{itemize}
      \pause
      \textbf{Gabarito:} (a)
    \item A principal justificativa para o uso de \emph{beamforming} em mmWave é:
      \begin{itemize}
        \item (a) reduzir o consumo de energia do UE
        \item (b) compensar a alta atenuação das ondas milimétricas
        \item (c) eliminar a numerologia OFDM
      \end{itemize}
      \pause
      \textbf{Gabarito:} (b)
    \item O \emph{network slicing} permite:
      \begin{itemize}
        \item (a) alocar recursos dedicados a eMBB, URLLC e mMTC sobre a mesma infraestrutura
        \item (b) aumentar a potência de transmissão do gNB
        \item (c) substituir o RAN por satélites LEO
      \end{itemize}
      \pause
      \textbf{Gabarito:} (a)
  \end{enumerate}
\end{frame}

\begin{frame}{Laboratório sugerido: 5G NR no ns-3}
  \begin{block}{Opção 1 --- Simulação de 5G NR}
    \begin{itemize}
      \item Usar o módulo \texttt{mmwave} (NYU) ou \texttt{nr} (5G-LENA) do ns-3: célula gNB--UE com \emph{beamforming} em 28~GHz.
      \item Medir \emph{throughput}, latência e SNR com UEs em mobilidade; comparar com um cenário sub-6~GHz.
    \end{itemize}
  \end{block}
  \begin{block}{Opção 2 --- Medição em rede móvel real}
    \begin{itemize}
      \item Medir a latência da rede móvel com \texttt{adb} (\emph{ping} contínuo) e \texttt{iperf3} entre smartphone e servidor.
      \item Estimar as parcelas de transmissão e propagação do RTT observado em 4G/5G.
    \end{itemize}
  \end{block}
\end{frame}

\section{Bibliografia}

\begin{frame}{Bibliografia}
  \begin{itemize}
    \item 3GPP TR 38.300, \emph{NR; Overall description; Stage-2}
    \item 3GPP TR 38.811, \emph{Study on New Radio (NR) to support non-terrestrial networks}
    \item 3GPP TR 38.821, \emph{Solutions for NR to support non-terrestrial networks (NTN)}
    \item 3GPP TS 23.501, \emph{System Architecture for the 5G System (5GS)}
    \item ITU-R M.2160, \emph{IMT Vision --- Framework and overall objectives of the future development of IMT for 2030 and beyond}
    \item M. Shafi et al., ``5G: A tutorial overview of standards, trials, challenges, deployment and practice'', \emph{IEEE JSAC} 35(6):1201--1221, 2017
  \end{itemize}
\end{frame}

\begin{frame}{Leitura recomendada}
  \begin{itemize}
    \item 3GPP TR 38.300, \emph{NR; Overall description; Stage-2} --- arquitetura geral e protocolos do 5G NR.
    \item 3GPP TS 23.501, \emph{System Architecture for the 5G System (5GS)} --- SBA, funções de rede e \emph{network slicing}.
    \item ITU-R M.2160, \emph{IMT Vision --- Framework and overall objectives of the future development of IMT for 2030 and beyond} --- requisitos e casos de uso das redes pós-5G.
    \item M. Shafi et al., ``5G: A tutorial overview of standards, trials, challenges, deployment and practice'', \emph{IEEE JSAC} 35(6):1201--1221, 2017.
  \end{itemize}
\end{frame}

\end{document}